\section{Interpretación conceptual de NIIF 17}
\label{sec:niif17}
\newcommand{\NPVbaseF}{-12\,912,86}
\newcommand{\MargenbaseF}{-57,88}
\newcommand{\RecursosUnoF}{1\,972,52}
\newcommand{\ReservaUnoF}{3\,997,13}
\newcommand{\BrechaUnoF}{-2\,024,61}


\textbf{1. Flujos iniciales y ajuste por riesgo.}
Para el seguro de C consideraríamos, dentro de los límites contractuales,
las primas esperadas de $50\,000$ por vigente, los beneficios por muerte de
$5\,000\,000$, los gastos directamente atribuibles de $10\,000$ por vigente,
la comisión del $5\%$ y los flujos de adquisición pertinentes de $75\,000$
por contrato. La proyección reflejaría fallecimientos, renuncias, momentos
de pago y descuento adecuado. El ajuste por riesgo no financiero (RA)
representa la compensación exigida por la incertidumbre en importes y
momentos de los flujos por riesgos como mortalidad, renuncia y gastos.
Esencialmente nos permite medir la asignación monetaria que se debe 
consignar asociado a la incertidumbre entre haber tenido un gasto fijo y otro variable.


\textbf{2. CSM y tratamiento inicial.}
El margen de servicio contractual (CSM) representa ganancia aún no devengada,
que se reconoce conforme se prestan servicios de seguro. En un grupo
rentable, el beneficio inicial pertinente se difiere mediante la CSM.
En un grupo oneroso, la pérdida inicial se reconoce de inmediato y la CSM
inicial es cero. El NPV negativo del ejercicio motiva revisar la rentabilidad,
pero la clasificación contable exige la medición y agrupación de NIIF 17, a saber por grupos de riesgo
 y rentabilidad compartida y también por banda temporal. Similarmente la CSM se debe de ajustar por periodos 
 considerando las unidades corrientes y las unidades de desarrollo futuro.

\textbf{3. Diferencia respecto del NPV.}
Un NPV positivo de \emph{profit testing} no demostraría por sí mismo una CSM
positiva: depende de las convenciones del ejercicio, a saber el descuento, el nivel de rendimiento, 
adquisición inicial y reservas suministradas. NIIF 17 mide grupos de contratos con flujos de cumplimiento, descuento y RA,
y considera sus límites contractuales y criterios de reconocimiento.
En nuestro caso, el NPV es \NPVbaseF{} millones de unidades monetarias y el
margen es \MargenbaseF{}\%; esos indicadores de rentabilidad no son una CSM
ni permiten cuantificarla. Por otro lado recordenmos que el NPV nos mide la rentabilidad esperada asociada a un contrato particular, el CSM considera la 
agrupación por contratos con características de rentabilidad , riesgo y temporalidad comúnes.

\textbf{4. CSM, reserva y efectivo.}
La CSM corresponde al beneficio por servicios futuros. La reserva académica
$V_t=v_tN_t$ sigue un cuadro fijo por vigente y no equivale automáticamente
al pasivo de NIIF 17. El efectivo corresponde a los recursos disponibles
para pagar obligaciones. Al cierre del año 1, antes de aportes, los recursos
son \RecursosUnoF{} millones y la reserva \ReservaUnoF{} millones, con una
diferencia de \BrechaUnoF{} millones. Por tanto, tener efectivo positivo
no asegura cubrir la reserva. La reserva no genera efectivo, su liberación
no es un cobro y un aporte de capital no constituye utilidad sino un aporte de capital que puede ayudar 
a cubrir las obligaciones futuras contraídas.

